\chapter{Coughing Up Blood}
\index{Coughing Up Blood}

\section*{Definition and Overview}
Coughing up blood (haemoptysis) is the spitting or coughing of blood from the airways or lungs. It ranges from small streaks of blood in the sputum to the coughing of large volumes of blood. Even a small amount of blood in the sputum is a significant symptom that requires medical assessment, because it can be the first sign of serious lung disease. Causes include bronchitis, infection, and, less commonly, lung cancer, pulmonary embolism, or bleeding disorders. The severity depends not only on the volume of blood but on the underlying cause, and large haemoptysis is a medical emergency.

\section*{Epidemiology}
Haemoptysis is a common reason for medical consultation. Most cases are minor and caused by respiratory infections or bronchitis, and the majority of people evaluated for haemoptysis are found to have benign causes. However, in smokers and older adults, haemoptysis can be the presenting feature of lung cancer, and in younger people it can indicate tuberculosis or other infections. Massive haemoptysis, though rare, is life-threatening, with a high mortality if not treated urgently.

\section*{Aetiology and Pathophysiology}
Blood in the sputum can come from any part of the respiratory tract.

\begin{itemize}
    \item \textbf{Bronchitis and infection:} inflammation of the airways, pneumonia, or tuberculosis
    \item \textbf{Lung cancer:} a tumour eroding into an airway or blood vessel
    \item \textbf{Bronchiectasis:} permanently widened, damaged airways that bleed easily
    \item \textbf{Pulmonary embolism:} a blood clot in the lung can cause bleeding and infarction
    \item \textbf{Heart failure:} raised pressure in the lung vessels can cause blood-stained sputum
    \item \textbf{Anticoagulant medicines:} blood thinners increase bleeding risk
    \item \textbf{Bleeding disorders:} low platelets or clotting problems
    \item \textbf{Trauma:} injury to the chest or airways
    \item \textbf{Rare causes:} autoimmune disease affecting the lungs, such as vasculitis or Goodpasture syndrome
\end{itemize}

\section*{Clinical Features}
The presentation depends on the volume and source of the bleeding.

\begin{itemize}
    \item Streaks or flecks of blood in otherwise clear or discoloured sputum
    \item Coughing up larger amounts of bright red or frothy blood
    \item Associated cough, fever, or chest pain
    \item Shortness of breath
    \item Weight loss, night sweats, or fatigue in chronic disease
    \item Features of anaemia with larger or repeated bleeds: pallor, dizziness, breathlessness
    \item A smoking history, which raises the possibility of lung cancer
\end{itemize}

\section*{Differential Diagnosis}
It is important to confirm that the blood is truly coming from the lungs.

\begin{itemize}
    \item Nosebleeds or bleeding from the throat and mouth dripping into the sputum
    \item Vomiting blood (haematemesis) from the stomach or oesophagus
    \item Bronchitis, pneumonia, and other respiratory infections
    \item Lung cancer and other tumours of the chest
    \item Pulmonary embolism
    \item Tuberculosis
    \item Heart failure and mitral stenosis
    \item Bleeding disorders and anticoagulant use
\end{itemize}

\section*{When to Seek Medical Care}
Anyone who coughs up blood should be assessed by a doctor promptly, even if the amount is small. Investigations are needed to establish the cause.

\textbf{Contact your local emergency services for:}
\begin{itemize}
    \item Coughing up a large amount of blood, or blood that continues or increases
    \item Difficulty breathing or severe shortness of breath
    \item Chest pain, dizziness, or collapse
    \item Coughing blood after a recent chest injury
    \item Signs of a serious bleed: pallor, rapid heart rate, or confusion
\end{itemize}

\section*{Diagnosis and Investigations}
The urgency and extent of investigation depend on the volume of bleeding and risk factors.

\begin{itemize}
    \item \textbf{History:} amount and colour of blood, smoking history, medicines, and associated symptoms
    \item \textbf{Examination:} chest examination, temperature, and oxygen saturation
    \item \textbf{Chest X-ray:} the first imaging test, looking for masses, infection, or fluid
    \item \textbf{Blood tests:} full blood count, clotting studies, inflammatory markers, and blood group
    \item \textbf{Sputum tests:} microscopy and culture for infection, including tuberculosis
    \item \textbf{CT scan:} detailed imaging of the lungs
    \item \textbf{Bronchoscopy:} a camera examination of the airways to find the bleeding site
    \item \textbf{Hospital admission:} for larger bleeds, for monitoring and urgent intervention
\end{itemize}

\section*{Management}
Treatment targets the underlying cause and, in severe cases, stops the bleeding.

\begin{itemize}
    \item \textbf{Mild haemoptysis from infection:} antibiotics for pneumonia or bronchitis and treatment of the infection
    \item \textbf{Antituberculous therapy:} for tuberculosis
    \item \textbf{Stopping or adjusting anticoagulants:} under medical supervision
    \item \textbf{Bronchoscopy:} to identify and sometimes treat the bleeding site
    \item \textbf{Interventional radiology:} embolisation of bleeding vessels
    \item \textbf{Surgery:} removal of a bleeding tumour or damaged lung segment in selected cases
    \item \textbf{Supportive care:} oxygen, fluids, and monitoring in hospital
    \item \textbf{Treatment of the underlying condition:} heart failure, bleeding disorders, or autoimmune disease
\end{itemize}

\section*{Complications}
\begin{itemize}
    \item Massive haemoptysis with airway obstruction and asphyxiation, a medical emergency
    \item Anaemia from blood loss
    \item Aspiration of blood into the lungs causing pneumonia
    \item Missed or delayed diagnosis of lung cancer or tuberculosis
    \item Complications of bronchoscopy or embolisation
    \item The underlying disease progressing if untreated
\end{itemize}

\section*{Prognosis and Outcomes}
The outlook depends entirely on the cause. Haemoptysis from acute bronchitis or infection usually settles as the infection resolves, and the prognosis is excellent. Haemoptysis from lung cancer, tuberculosis, or bronchiectasis is a symptom of an underlying condition that needs treatment, and the outcome follows the prognosis of that disease. Massive haemoptysis carries a high mortality and requires urgent specialist care. Because early diagnosis of serious disease improves outcomes, prompt investigation of even small bleeds is important.

\section*{Prevention and Self-Care}
\begin{itemize}
    \item Do not smoke; smoking is the leading cause of serious lung disease
    \item Seek medical assessment for any blood in the sputum, even a small amount
    \item Keep a sample of the blood-stained sputum to show the doctor
    \item Take anticoagulant medicines exactly as prescribed and have blood tests as directed
    \item Treat chest infections promptly, especially in people with chronic lung disease
    \item Attend follow-up so that the cause is confirmed and treated
    \item Seek urgent care for any increase in bleeding or difficulty breathing
\end{itemize}

\section*{Sources and Further Reading}
The following reputable sources provide further information for healthcare students and patients seeking reliable, current advice about coughing up blood.

\begin{itemize}
    \item NHS. \textit{Coughing up blood (haemoptysis).} https://www.nhs.uk/conditions/coughing-up-blood/
    \item Mayo Clinic. \textit{Coughing up blood: causes and when to see a doctor.} https://www.mayoclinic.org/symptoms/coughing-up-blood/
    \item MedlinePlus. \textit{Haemoptysis (coughing up blood).} https://medlineplus.gov/ency/article/003073.htm
    \item MSD Manual. \textit{Haemoptysis: professional reference.} https://www.msdmanuals.com/
    \item Healthdirect Australia. \textit{Coughing up blood.} https://www.healthdirect.gov.au/
    \item Lung Foundation Australia. \textit{Information on lung conditions and symptoms.} https://lungfoundation.com.au/
\end{itemize}
